% 39-05(39쪽) — 구간 [0,4]에서 정의된 y=f(x)의 그래프(올라간 뒤 꺾여 내려오다 x=3 에서 끊기고 다시 내려감 · 빈 점과 채운 점). 스캔 화질이 낮아 TikZ 로 다시 그림.
% 원문에 있는 것만: 채운 점 (0,1), (3,2), (4,2) · 빈 점 (3,3), (4,0) · 안내 점선 · 눈금 1, 2, 3 과 1, 2, 3, 4 · 라벨 O, x, y, y=f(x).
\begin{tikzpicture}[x=0.45cm, y=0.45cm, line width=0.5pt, font=\small]
  \draw[->, >=stealth, line width=0.7pt] (-0.5,0) -- (4.85,0) node[below=1pt] {$x$};
  \draw[->, >=stealth, line width=0.7pt] (0,-0.7) -- (0,3.85) node[left=1pt] {$y$};
  \draw[densely dashed, line width=0.4pt] (0,3) -- (3,3);
  \draw[densely dashed, line width=0.4pt] (0,2) -- (4,2);
  \draw[densely dashed, line width=0.4pt] (1,0) -- (1,2);
  \draw[densely dashed, line width=0.4pt] (2,0) -- (2,3);
  \draw[densely dashed, line width=0.4pt] (3,0) -- (3,3);
  \draw[densely dashed, line width=0.4pt] (4,0) -- (4,2);
  \draw[line width=0.6pt] (0,1) -- (2,3) -- (3,2);
  \draw[line width=0.6pt] (3,3) -- (4,0);
  \fill (0,1) circle (2.2pt);
  \fill (3,2) circle (2.2pt);
  \fill (4,2) circle (2.2pt);
  \draw[fill=white, line width=0.5pt] (3,3) circle (2.2pt);
  \draw[fill=white, line width=0.5pt] (4,0) circle (2.2pt);
  \node[left=1pt] at (0,1) {$1$};
  \node[left=1pt] at (0,2) {$2$};
  \node[left=1pt] at (0,3) {$3$};
  \node[below left=-2pt and -3pt] at (0,0) {$\pt{O}$};
  \node[below=1pt] at (1,0) {$1$};
  \node[below=1pt] at (2,0) {$2$};
  \node[below=1pt] at (3,0) {$3$};
  \node[below=1pt] at (4,0) {$4$};
  \node[anchor=west] at (1.6,3.6) {$y=f(x)$};
\end{tikzpicture}
